\section{Normalized volume and the cone reduction}\label{sec:cones}

\subsection{Conventions and public inputs}

All varieties and stacks in the proof are over \(\CC\). Boundaries
are zero unless a boundary is explicitly introduced in a vanishing
argument. For a klt germ \(x\in X\), write \(R=\cO_{X,x}\) and
\(\fm=\fm_x\). If \(E\) is a prime divisor on a smooth proper
birational model \(\mu\colon X'\to X\), its log discrepancy is
\[
 A_X(E)=1+\operatorname{coeff}_E(K_{X'}-\mu^*K_X).
\]
For a real valuation \(v\) of \(K(X)\), trivial on \(\CC\) and
centered exactly at \(x\), put
\[
 \fa_t(v)=\{f\in R:v(f)\ge t\},\qquad
 \vol(v)=\lim_{t\to\infty}
     \frac{n!\,\length(R/\fa_t(v))}{t^n}.
\]
The ideals are \(\fm\)-primary. The limit exists by
\cite[Theorem~5.8]{Cutkosky13}: the local rings here are analytically
unramified and equicharacteristic with perfect residue field. The
integer-cutoff ideals form a graded family, and monotonicity extends
the integer limit to real cutoffs.

We use the standard extension of log discrepancy to real valuations.
On a smooth birational log model it is linear on each monomial cone;
for an arbitrary valuation it is the supremum over its monomial
retractions \cite[Proposition~5.1 and Corollary~5.4]{JM12}.
For the singular \(\QQ\)-Gorenstein base one includes the relative
canonical correction on a resolution, as in
\cite[Section~2.1]{Li18}. Define
\[
 \nvol_X(v)=
 \begin{cases}
 A_X(v)^n\vol(v),&A_X(v)<+\infty,\\
 +\infty,&A_X(v)=+\infty.
 \end{cases}
\]
Thus the definition never requires multiplying infinity by zero.
Our normalization is
\begin{equation}\label{cone:eq-density}
 \nvol(x,X)=\inf_v\nvol_X(v),\qquad
 d(x,X)=\frac{\nvol(x,X)}{n^n}.
\end{equation}
The infimum is unchanged if restricted to positive multiples of
divisorial valuations \cite[Theorem~2.5]{LiXu18}.
Rescaling a valuation does not change its
normalized volume.

We use the following established facts in their boundary-zero form.
\begin{enumerate}[label=(\roman*)]
\item A klt germ has a minimizing valuation, unique up to positive
scaling. It has finite discrepancy, is quasi-monomial, and has
finitely generated associated graded. The latter is a K-semistable
Fano cone, which admits a special equivariant degeneration to a
K-polystable Fano cone. The minimizing Reeb vector and normalized
volume are preserved \cite[Theorem~1.2]{XZStable},
\cite[Theorem~1.2]{LWX21}. Uniqueness also implies that a minimizer
normalized by discrepancy is invariant under every automorphism of
the germ \cite[Theorem~1.1 and Corollary~1.2]{XZ21}.
\item Normalized volume is lower semicontinuous in
\(\QQ\)-Gorenstein flat families of klt germs with section over a normal base
\cite[Theorem~1]{BL21}. In particular it is lower semicontinuous
when the point moves along a curve in a fixed klt variety
\cite[Theorem~33(1)]{BL21}. One has
\(d(x,X)\le1\), with equality exactly at smooth points
\cite[Theorem~1.6]{LX19}.
\item For a finite Galois morphism of klt germs
\(f\colon (Y,y)\to(X,x)\), with \(f^{-1}(x)=\{y\}\) set-theoretically
and \(K_Y=f^*K_X\), the finite degree formula is
\begin{equation}\label{cone:eq-cover}
              \nvol(y,Y)=\deg(f)\nvol(x,X).
\end{equation}
We use this only for a quotient by a group without reflections and
for the cyclic canonical index cover
\cite[Theorem~1.3]{XZ21}.
\item Isomorphic completed local rings have the same normalized
volume \cite[Lemma~2.15]{Liu22}. Thus étale germs at complex closed
points have the same invariant.
\end{enumerate}
These facts concern arbitrary klt germs. In particular, the first
item does not require a divisorial minimizer, isolatedness, or
smoothability.

\subsection{Products and embedding dimension}

\begin{lemma}[A smooth factor]\label{cone:product}
Let \(y\in Y\) be an \(N\)-dimensional klt germ, where \(N\ge1\). Then
\[
          d((y,0),Y\times\A^1)\le d(y,Y).
\]
\end{lemma}

\begin{proof}
Take a divisorial test \(v\) centered at \(y\), of discrepancy
\(A>0\) and volume \(\sigma\), and let \(a=A/N\).
The Gauss extension to the added parameter \(t\) is
\[
       (v\oplus a)\left(\sum_i f_it^i\right)
                    =\min_i\{v(f_i)+ia\}.
\]
On a log resolution this is a monomial valuation along the divisor
of \(v\) and \(t=0\), so its log discrepancy is \(A+a\).
Completion does not change the finite colengths. In the power-series
ring the quotient decomposes by powers of \(t\), giving
\[
 \length\bigl(R[[t]]/\fa_k(v\oplus a)\bigr)
   =\sum_{0\le i<k/a}\length\bigl(R/\fa_{k-ia}(v)\bigr).
\]
The defining asymptotic for \(\sigma\), followed by a Riemann sum,
gives \(\vol(v\oplus a)=\sigma/a\). Therefore
\[
 \frac{(A+A/N)^{N+1}}{(N+1)^{N+1}}\frac{N\sigma}{A}
                   =\frac{A^N\sigma}{N^N}.
\]
Take the infimum over the divisorial tests on \(Y\).
\end{proof}

\begin{lemma}[Lifting generators]\label{cone:edim}
Let \(v\) be a finite-discrepancy valuation of a klt germ, with
finitely generated positively graded associated graded ring
\(\gr_v R\). Then
\[
                \edim(R)\le\edim(\gr_v R).
\]
\end{lemma}

\begin{proof}
Choose minimal homogeneous algebra generators
\(\bar f_1,\ldots,\bar f_e\) of \(\gr_v R\), all of positive value,
and lift them to \(f_1,\ldots,f_e\in\fm\). To show that their
classes span \(\fm/\fm^2\), start with \(g\in\fm\). Its initial
form is a homogeneous polynomial in the \(\bar f_i\). Subtracting
the corresponding polynomial in the \(f_i\) either annihilates
the remainder or strictly increases its value. Every positive
value that occurs belongs to the semigroup generated by the
positive numbers \(v(f_i)\), so only finitely many such values
lie below any fixed bound.

The Izumi inequality for a finite-discrepancy valuation on a klt
germ \cite[Theorem~3.1]{Li18} gives
\(v(g)\le c\,\ord_{\fm}(g)\), with a fixed constant
\(c\). In particular, a nonzero remainder of value greater than
\(c\) lies in \(\fm^2\). The cancellation process therefore
terminates modulo \(\fm^2\) after finitely many steps. Every
polynomial subtracted has no constant term, and modulo
\(\fm^2\) is a linear combination of the \(f_i\). This proves
the assertion.
\end{proof}

\subsection{The minimum at the ordinary double point}

\begin{lemma}[The ODP minimum]\label{end:odp-value}
For \(n\ge2\), the ordinary double point \(Q_n\) has
\[
                           \nvol(0,Q_n)=2(n-1)^n.
\]
\end{lemma}

\begin{proof}
First take the exact quadratic cone
\[
 Q_n=\Spec R,\qquad
 R=\CC[z_0,\ldots,z_n]/(z_0^2+\cdots+z_n^2),
\]
with its ordinary grading. It is klt. For example, blowing
up the vertex resolves it with smooth exceptional quadric,
whose log discrepancy is \(n-1>0\); further smooth
blowups preserve the klt condition.

Normalize a minimizing valuation \(v\) by its discrepancy.
Uniqueness makes it invariant under all automorphisms,
in particular under cone scalings and the connected
orthogonal group \(\operatorname{SO}(n+1,\CC)\).
For each \(j\ge0\), the degree-\(j\) piece \(R_j\) is
the irreducible representation of complex harmonic polynomials
of degree \(j\) in \(n+1\) variables. To identify the quotient,
write \(P_j\) for homogeneous complex polynomials of degree \(j\)
and \(\mathcal H_j\) for the subspace killed by
\(\sum_i\partial_{z_i}^2\); set \(P_k=0\) for \(k<0\).
For \(q=\sum z_i^2\), the harmonic decomposition
\(P_j=\mathcal H_j\oplus qP_{j-2}\) identifies \(R_j\) with
\(\mathcal H_j\); see
\cite[Example~251 and Theorem~257]{GrahamSwansonMatthews15}.
The complex representation is irreducible already under the compact
group \(\operatorname{SO}(n+1)\), so an invariant subspace under
\(\operatorname{SO}(n+1,\CC)\) must be trivial or the whole space.
For the classical harmonic-polynomial formulation and its extensions,
see also \cite[Section~1]{LuoXu12}.
This includes \(n=2\) and \(n=3\).
The exception for two ambient variables is outside
the range under consideration.

For each cutoff \(t\), the subspace
\(\{f\in R_j:v(f)\ge t\}\) is invariant, so it is
either zero or all of \(R_j\). Every nonzero element of
\(R_j\) consequently has a common value \(a_j\).
If \(0\ne f\in R_1\), then \(f^j\ne0\), and
valuation multiplicativity gives \(a_j=j a_1\).
Invariance under cone scalings separates homogeneous
pieces of a cutoff: a finite Vandermonde combination
of scalings isolates each piece. Thus, for
\(g=\sum_j g_j\ne0\),
\[
                       v(g)=a_1\min\{j:g_j\ne0\}.
\]
After localizing at the vertex the same assertion
holds, since denominators nonzero at the vertex have
value zero. The minimizer is therefore proportional
to the ordinary degree valuation.

That valuation is the order along the exceptional
quadric of the vertex blowup. Hypersurface adjunction
gives discrepancy \(n-1\), and its volume is the
hypersurface multiplicity \(2\). It follows that the
infimum equals \(2(n-1)^n\).

Finally, an analytic ordinary double point has the
same completed local ring as this exact cone.
Completion invariance \cite[Lemma~2.15]{Liu22}
gives the same normalized volume.
\end{proof}

\subsection{The induction setup}

The companion paper \emph{The normalized-volume gap in dimension four}
\cite[Theorem~1.1]{FourfoldCompanion} proves the following base case.
We state the entire interface because its equality assertion is part
of the theorem propagated here.

\begin{theorem}[The fourfold gap]\label{inp:fourfold}
Let \(x\) be a singular closed point of a normal complex algebraic
fourfold \(X\). Suppose that \(K_X\) is \(\QQ\)-Cartier and \(X\) is klt
near \(x\), with zero boundary. Then
\[
                  \nvol(x,X)\le162.
\]
Equality holds if and only if the complex analytic germ is
\[
       \left(\{z_0^2+z_1^2+z_2^2+z_3^2+z_4^2=0\},0\right).
\]
No isolatedness, complete-intersection, quotient, or smoothability
assumption is made.
\end{theorem}

For \(k\ge2\) set
\begin{equation}\label{cone:eq-threshold}
              p_k=2\left(\frac{k-1}{k}\right)^k.
\end{equation}
Let \(M_k\) be the supremum of \(d(x,X)\) over all singular
boundary-zero klt \(k\)-fold germs: the varieties are normal and
complex algebraic, and their canonical divisors are
\(\QQ\)-Cartier. This supremum is finite, since \(M_k\le1\).
The numbers \(p_k\) are strictly increasing: for real \(t>1\),
\[
 \frac{d}{dt}\left[t\log(1-1/t)\right]
   =\log(1-1/t)+\frac1{t-1}>0.
\]
The last inequality is \(\log(1+u)<u\), with \(u=1/(t-1)\).

The dimension-two case of Theorem~\ref{thm:main} follows from
the quotient classification of klt surface singularities
\cite[Proposition~6.11]{CKM88}. The germ is analytically isomorphic to
\((\A^2/G,0)\), where \(G\) is a finite effective linear group
acting without reflections. Analytic isomorphism identifies the
completed local rings, so completion invariance transfers the
normalized volume to this algebraic quotient. The finite
degree formula gives \(d=1/|G|\). If the germ is singular,
\(|G|\ge2\), and equality with \(p_2=1/2\) forces an order-two
action without a fixed hyperplane, namely \((-1,-1)\).
This is the \(A_1\) surface singularity. The threefold case is
\cite[Theorem~1.3(1)]{LX19}, and the fourfold case is
Theorem~\ref{inp:fourfold}.

Henceforth, until the induction is closed, assume
Theorem~\ref{thm:main} in dimensions below \(n\), and fix
\begin{equation}\label{cone:eq-induction}
 n=N+1\ge5,\qquad
 d(x,X)\ge p_n,
\end{equation}
with \(x\) singular. The only use of \(M_n\) will be its definition
as a supremum, not an unproved upper bound in dimension \(n\).

\begin{proposition}[Cone reduction]\label{cone:reduction}
Under \eqref{cone:eq-induction}, there is a singular
K-polystable Fano cone \(C=\Spec S\), with vertex \(o\), such that
\[
 d(o,C)=d(x,X),\qquad
 \edim(X,x)\le\edim(C,o).
\]
The puncture \(C\setminus\{o\}\) is smooth, \(S\) is a normal
Gorenstein ring, and its canonical module has a homogeneous
generator. The cone has zero boundary.
\end{proposition}

\begin{proof}
Apply stable degeneration to a minimizer at \(x\), and then take
the K-polystable special degeneration. In the first passage the
induced Reeb valuation has the same volume and discrepancy as the
original minimizer \cite[Lemma~4.10]{LiXu18}.
In the second, equivariant flatness preserves the dimensions of
the character spaces and hence the Reeb volume. The relative
canonical character is constant, so its discrepancy is also
preserved. Li--Xu's characterization of K-semistability
\cite[Theorem~1.3]{LiXu18} shows that the final Reeb valuation
minimizes normalized volume over all centered valuations. Thus the two passages
preserve the infimum in \eqref{cone:eq-density}, not only the
number of an arbitrary test.

Lemma~\ref{cone:edim} gives the embedding-dimension inequality
for the first passage. Upper semicontinuity along the vertex
section of the special flat degeneration gives it for the
second. The final vertex is therefore singular.

Choose any positive integral grading of \(S\), and let \(P\ne o\).
Its grading orbit specializes to \(o\), so lower semicontinuity
gives
\begin{equation}\label{cone:eq-orbit-density}
                         d(o,C)\le d(P,C).
\end{equation}
Choose a homogeneous function \(f\), of degree \(d>0\), with
\(c=f(P)\ne0\), and put \(Y_P=\{f=c\}\subset C_f\).
The action map
\[
 Y_P\times\Gm\longrightarrow C_f,\qquad (Q,a)\longmapsto a\cdot Q,
\]
is the finite étale cover obtained by adjoining a root
\(a^d=f/c\). Near \((P,1)\), it expresses the germ as the product
of the \(N\)-dimensional slice \((Y_P,P)\) and a smooth curve.
The slice is klt by this étale description.
If \(P\) were singular, the slice would be
singular as well. The induction hypothesis,
Lemma~\ref{cone:product}, and étale invariance would then imply
\[
              d(P,C)\le p_N<p_n\le d(o,C),
\]
contrary to \eqref{cone:eq-orbit-density}. Thus the puncture is
smooth.

Let \(q\) be the canonical index at \(o\). The cyclic index
cover has degree \(q\), is quasi-étale and crepant, and is klt.
It has one point over \(o\). Indeed, its algebra is the direct
sum of the reflexive canonical powers modulo \(q\).
If a homogeneous element of a nonzero residue degree were
not nilpotent in the fiber over \(o\), a power returning to
degree zero would be a unit. The corresponding section would
trivialize a proper positive canonical power, contradicting
minimality of \(q\). Thus the nonzero residue-degree pieces
are nilpotent, and the reduced fiber consists of one point.
Equation~\eqref{cone:eq-cover} and the upper bound \(d\le1\)
give \(d(o,C)\le1/q\). Since \(p_n>1/2\), necessarily \(q=1\).

Klt singularities are Cohen--Macaulay
\cite[Definition~117 and Theorem~120]{KollarCanonicalModels10}.
The canonical module
\(\omega_S\) is invertible at the vertex, and a homogeneous
element generating its one-dimensional fiber there generates
it after localization. Its graded cokernel has empty support:
any nonempty closed support invariant under a positive
grading contains the vertex. Since \(\omega_S\) is torsion-free,
the resulting map \(S\to\omega_S\) is injective as well.
It is therefore an isomorphism, with a grading shift.
\end{proof}

\subsection{Integral gradings and orbifold quotients}

Fix the cone of Proposition~\ref{cone:reduction}. Let \(\mathbb T\)
be an effective algebraic torus of automorphisms containing its
Reeb vector \(\xi\), and normalize
\[
                              A(\xi)=n.
\]
Choose a \(\mathbb T\)-eigen-generator \(\Omega_C\) of
\(\omega_S\). Its character pairs with cocharacters as the
linear discrepancy functional \(A\)
\cite[Lemmas~2.16 and~2.18]{LiXu18}. Function weights use the
convention \(g(aP)=a^j g(P)\); the differential \(dg\) has the
same weight. The Reeb cone consists of cocharacters positive
on all nonconstant homogeneous functions.

\begin{lemma}[Fast primitive approximations]\label{cone:approximation}
There are primitive positive integral cocharacters \(\zeta_i\),
of canonical weights \(r_i=A(\zeta_i)\in\ZZ_{>0}\), such that
\begin{equation}\label{cone:eq-fast}
             \zeta_i=\frac{r_i}{n}\xi+\epsilon_i,\qquad
             \epsilon_i\longrightarrow0
\end{equation}
in the real cocharacter space. If the ray of \(\xi\) is
rational, the sequence may be chosen constant and exact.
Otherwise \(r_i\to\infty\).

The normalized volumes of the degree valuations
\(\wt_{\zeta_i}\) tend to \(\nvol(o,C)\).
For any fixed finite set of homogeneous algebra generators,
their \(\zeta_i\)-degrees lie between \(cr_i\) and \(Cr_i\)
for positive constants \(c,C\). Consequently the largest
stabilizer order \(m_{\max,i}\) on the puncture is \(O(r_i)\).
\end{lemma}

\begin{proof}
In the quotient of real cocharacter space by its integral
lattice, recurrence gives \(t_i\to\infty\) and integral
cocharacters \(q_i\) with \(q_i=t_i\xi+e_i\), where \(e_i\to0\).
One may obtain such a sequence from simultaneous Dirichlet
approximation, or from returns to zero of the one-parameter
subgroup in its compact closure. Let \(d_i\) be the gcd of
\(q_i\) in a lattice basis and put \(\zeta_i=q_i/d_i\).
Then
\[
 \zeta_i-\frac{A(\zeta_i)}n\xi
       =\frac1{d_i}\left(e_i-\frac{A(e_i)}n\xi\right)\longrightarrow0.
\]
The cocharacters are positive for large \(i\).
If their positive integers \(r_i=A(\zeta_i)\) had a bounded
subsequence, \eqref{cone:eq-fast} would make \(\zeta_i\)
bounded on that subsequence. A constant lattice subsequence
would then give \(\zeta=(r/n)\xi\), forcing the ray to be
rational. The rational case is handled by its primitive
generator.

Choose homogeneous algebra generators \(g_1,\ldots,g_a\)
with \(\mathbb T\)-characters \(\chi_1,\ldots,\chi_a\).
Each \(\langle\chi_j,\xi\rangle\) is positive, so
\eqref{cone:eq-fast} gives the asserted two-sided bounds
for their integral degrees. More precisely, the normalized
cocharacters \((n/r_i)\zeta_i\) converge to \(\xi\).
Uniform comparison on the finite set of generator characters
then compares their values on every occurring character:
each is a sum of generator characters with nonnegative
integer coefficients. It follows that the valuation ideals
of the normalized degree valuations are squeezed between
those of \(\wt_\xi\) with cutoff factors tending to one.
Their volumes therefore converge; their discrepancies are
all \(n\).

At a nonvertex point some generator is nonzero. Its
stabilizer for a primitive grading is cyclic, and its order
divides the degree of every nonvanishing generator. It is
therefore at most the largest generator degree, proving
the final assertion.
\end{proof}

In what follows, \emph{sufficiently fast gradings} always means
gradings in a sequence as in Lemma~\ref{cone:approximation},
including the exact constant sequence on a rational ray.
All constants may depend on the fixed cone. No uniformity
over the class defining \(M_n\) is asserted or needed.

\begin{lemma}[The quotient orbifold]\label{cone:orbifold}
For a primitive positive integral grading \(\zeta\), of
canonical weight \(r\), let
\[
          B=[(C\setminus\{o\})/\Gm].
\]
Then \(B\) is a smooth effective proper orbifold of dimension
\(N\), with projective coarse space \(\Proj S\).
It has a faithful ample weight-one line bundle \(L\) such that
\begin{equation}\label{cone:eq-section-ring}
         S_j=H^0(B,jL),\qquad -K_B=rL.
\end{equation}
Every stabilizer is a finite cyclic group. At a point of
order \(m\), a smooth transverse uniformizer carries a
faithful \(\mu_m\)-action, and \(L\) has a faithful fiber
character.
\end{lemma}

\begin{proof}
The smooth puncture has finite stabilizers for the positive
grading, so its quotient stack is smooth and Deligne--Mumford.
It is the usual stack associated with a positively graded
affine ring away from its vertex; its coarse space is
\(\Proj S\), and a sufficiently divisible positive power of
the weight-one line descends to an ample coarse line bundle.
This also proves properness.

The grading is effective because the integral cocharacter
is primitive in an effective torus. At a point, the finite
stabilizer acts trivially on the tangent line of the orbit.
If an element acted trivially also on a transverse tangent
space, it would have identity differential on the smooth
puncture. A finite-order automorphism in characteristic zero
is formally linearizable at a fixed point; hence that
element would be the identity on a neighborhood and thus
on the irreducible cone, contrary to effectiveness.
This proves faithfulness of the transverse representation
and effectiveness of the orbifold. The weight-one character
of the subgroup of \(\Gm\) is faithful by construction.

Sections of \(jL\) are the homogeneous functions of degree
\(j\) on the puncture. Normality and the codimension of the
vertex extend them to \(C\), giving the first equality in
\eqref{cone:eq-section-ring}. Contracting \(\Omega_C\)
with the Euler field gives a horizontal top form of
character \(r\), so \(K_B=-rL\).
\end{proof}

We retain the stack \(B\), its faithful line \(L\), and the
character information throughout the argument. In particular,
we do not replace \(L\) by a coarse line bundle before the
stabilizers have been shown to be trivial.
